\documentclass[12pt,a4paper]{article}

% --- Packages --------------------------------------------------------------
\usepackage[a4paper, margin=1in]{geometry}
\usepackage{times}
\usepackage{setspace}
\usepackage{titlesec}
\usepackage{booktabs}
\usepackage{array}
\usepackage{tabularx}
\usepackage{longtable}
\usepackage{multirow}
\usepackage[table]{xcolor}
\usepackage{hyperref}
\usepackage{enumitem}
\usepackage{parskip}
\usepackage{microtype}
\usepackage{caption}
\usepackage{fancyhdr}
\usepackage{graphicx}
\usepackage{tikz}
\usetikzlibrary{shapes.geometric, arrows.meta, positioning, fit, backgrounds, calc}
\usepackage{amsmath}

% --- Colours ---------------------------------------------------------------
\definecolor{darkblue}{RGB}{31,56,100}
\definecolor{midblue}{RGB}{46,95,163}
\definecolor{lightblue}{RGB}{235,243,251}
\definecolor{tableheadbg}{RGB}{31,56,100}
\definecolor{altrow}{RGB}{245,249,255}
\definecolor{sensecol}{RGB}{200,225,255}
\definecolor{detectcol}{RGB}{200,240,210}
\definecolor{analysecol}{RGB}{255,230,180}
\definecolor{riskcol}{RGB}{255,200,200}
\definecolor{outputcol}{RGB}{220,200,255}
\definecolor{senseframe}{RGB}{31,90,160}
\definecolor{detectframe}{RGB}{30,130,60}
\definecolor{analyseframe}{RGB}{180,100,0}
\definecolor{riskframe}{RGB}{180,30,30}
\definecolor{outputframe}{RGB}{100,30,160}

% --- Hyperref --------------------------------------------------------------
\hypersetup{
  colorlinks=true,
  linkcolor=midblue,
  urlcolor=midblue,
  citecolor=midblue
}

% --- Section formatting ----------------------------------------------------
\titleformat{\section}{\large\bfseries\color{darkblue}}{{\thesection.}}{0.5em}{}[\vspace{-4pt}\rule{\linewidth}{0.6pt}\vspace{2pt}]
\titleformat{\subsection}{\normalsize\bfseries\color{midblue}}{\thesubsection}{0.5em}{}
\titleformat{\subsubsection}{\normalsize\bfseries\color{darkblue}}{\thesubsubsection}{0.5em}{}

% --- Header / Footer -------------------------------------------------------
\pagestyle{fancy}
\fancyhf{}
\rhead{\small\color{gray}Crowd Flow \& Density Safety Oracle}
\lhead{\small\color{gray}Project Report}
\rfoot{\small\thepage}
\renewcommand{\headrulewidth}{0.4pt}

% --- Spacing ---------------------------------------------------------------
\onehalfspacing
\setlength{\parindent}{0pt}
\setlength{\parskip}{6pt}

% --- Table column types ----------------------------------------------------
\newcolumntype{L}[1]{>{\raggedright\arraybackslash}p{#1}}
\newcolumntype{C}[1]{>{\centering\arraybackslash}p{#1}}

% ===========================================================================
\begin{document}

% --- Title Page ------------------------------------------------------------
\begin{titlepage}
  \centering
  \vspace*{2cm}
  {\Huge\bfseries\color{darkblue} Project Report\par}
  \vspace{0.5cm}
  \rule{\linewidth}{1.2pt}
  \vspace{0.5cm}
  {\Large\bfseries\color{midblue}
    Intelligent Crowd Flow \& Density Safety Oracle\\[0.3em]
    \normalsize\color{gray}(Smart Venues / Event Management)
  \par}
  \vspace{0.5cm}
  \rule{\linewidth}{1.2pt}
  \vspace{1.5cm}
  {\normalsize Submitted as part of Academic Project Report\par}
  \vspace{0.3cm}
  {\normalsize April 2026\par}
  \vfill
\end{titlepage}

\newpage
\tableofcontents
\newpage

% ===========================================================================
\section{Introduction}
% ===========================================================================

Crowd management at large-scale public venues has historically been a reactive discipline. Security personnel stationed at fixed points, manually operated turnstiles, and post-incident forensic analysis have constituted the standard safety toolkit for stadium events, religious congregations, music festivals, and transit hub peak hours alike. As urban populations grow and the frequency of mass gatherings increases, this reactive paradigm is proving fundamentally inadequate. The tragic consequences of unmanaged crowd dynamics — stampedes, crush events, and panic-induced injuries — continue to claim lives globally, despite widespread deployment of closed-circuit television infrastructure and the availability of digital ticketing data.

The emergence of edge artificial intelligence represents a transformative opportunity for crowd safety engineering. By deploying machine intelligence at the point of sensing — on affordable, low-power hardware mounted directly at venue locations — rather than routing data to centralised cloud servers, it becomes possible to greatly reduce network latency from the safety decision loop. When combined with real-time computer vision, directional flow analysis, and predictive risk scoring, edge AI enables a fundamental shift from reactive monitoring to proactive crowd governance.

This project presents the \textbf{Intelligent Crowd Flow \& Density Safety Oracle}: a system that fuses YOLOv8-tiny object detection, clustering-based density heatmaps, and optical flow vector analysis on a Raspberry Pi edge node paired with an ESP32-CAM camera module. Risk scores are published via MQTT to a central dashboard, which simultaneously drives automated entry and exit signal actuation through an ESP32 microcontroller. The current demonstration is implemented as a single-camera setup using one Raspberry Pi node, which successfully demonstrates the complete end-to-end pipeline. However, the architecture is inherently modular and designed to easily scale up to $n$ independent camera nodes deployed across multiple venue locations for real-world deployment.

This report is organised as follows. Section~2 presents the problem statement. Section~3 surveys the relevant literature. Section~4 provides a comparative analysis. Section~5 identifies research gaps. Section~6 establishes the novelty of the proposed approach. Section~7 describes the complete system methodology.

% ===========================================================================
\section{Problem Statement}
% ===========================================================================

Large-scale public gatherings represent one of the most complex and under-addressed challenges in public safety engineering. Incidents such as the Seoul Halloween crowd crush (2022, 159 fatalities), the Morbi Bridge collapse (2022, 135 fatalities), and the Karur Temple stampede (2024) have repeatedly demonstrated the catastrophic cost of reactive crowd management strategies. In each case, the contributing factors were identical: undetected density accumulation, unchecked unidirectional crowd flow, the absence of real-time risk quantification, and the complete absence of automated, closed-loop venue actuation to redirect or moderate crowd movement before a critical threshold was breached.

Existing technological approaches to crowd monitoring fall into three categories, each with fundamental limitations. Manual surveillance relies on human operators watching camera feeds — cognitively unreliable at scale and inherently reactive. Traditional server-side deep learning systems offer improved accuracy but depend on centralised infrastructure, introduce network latency incompatible with time-critical safety decisions, and produce counting outputs rather than actionable risk scores. Specialised thermal or radar sensor systems address illumination challenges but sacrifice the spatial and directional resolution needed for flow-vector analysis at crowd chokepoints.

Critically, no existing deployed or published system simultaneously achieves: (1) real-time crowd density and flow-vector computation on affordable, venue-mountable edge hardware without dependency on centralised cloud infrastructure; (2) proactive stampede and crush-risk prediction based on both density gradient and flow divergence trends rather than reactive threshold breach detection; and (3) closed-loop, autonomous actuation of physical venue entry and exit signals without human-in-the-loop delays. This triple gap — in deployment cost, proactivity, and actuation — defines the core engineering problem this project addresses.

% ===========================================================================
\section{Literature Survey}
% ===========================================================================

\subsection{Crowd Counting and Density Estimation}

Crowd density estimation has evolved from regression-based methods using hand-crafted features (HOG, LBP, SIFT) to CNN-based density map regression and transformer models. A 2025 Springer survey~\cite{survey2025} documents this progression, noting that detection-based approaches such as YOLO~\cite{yolo} and SSD variants fail in ultra-dense crowds due to severe occlusion. A 2024 IEEE paper~\cite{ieee_yolo_crowd} evaluated YOLOv3 with Canny Edge Detection (86.8\% accuracy) and YOLOv5 with OpenCV, demonstrating YOLO-family viability for crowd counting — though without flow dynamics, risk prediction, or edge deployment. More recent CNN-based approaches~\cite{nature_crowd_cnn,ieee_xai_crowd} advance server-side density estimation with explainability but remain cloud-reliant without any actuation output.

\subsection{Stampede and Crush Risk Detection}

An IEEE 2025 system~\cite{ieee_thermal} using thermal cameras computes crowd density in real time and transmits results for proactive safety, but lacks optical flow analysis and venue actuation. A 2024 paper from Beijing Jiaotong University~\cite{detr_stampede} introduced a Tiny-Scale Strengthened Deformable DETR stampede alert system achieving a 34\% improvement in small-target detection, yet operates reactively on GPU servers without venue integration. A 2023 ScienceDirect study~\cite{caem_2023} combined YOLOv3 with a Cellular Automata Evacuation Model at 96.5\% detection accuracy, generating Gaussian kernel heatmaps — but remains cloud-server dependent with no live actuation.

\subsection{Crowd Flow Analysis and Anomaly Detection}

Optical flow techniques — Lucas-Kanade and Farneback — have been applied to extract motion vector fields from crowd video. A 2025 arXiv study~\cite{farneback_2025} employed event-driven Farneback sampling to reduce computational load while maintaining sensitivity, benchmarked on ShanghaiTech. The active-Langevin model~\cite{langevin} advanced flow segmentation for pre-stampede collective motion. The 2025 mmFlux framework~\cite{mmflux} applied mmWave MIMO radar with optical flow principles to build semantic crowd inference graphs — weather-robust but lacking individual-level granularity and requiring dedicated infrastructure.

\subsection{Edge AI Deployment and IoT Middleware}

A 2026 Scientific Reports paper~\cite{yolov8_edge_2026} confirmed that YOLOv8-tiny is a viable model choice for Raspberry Pi platforms running lightweight inference pipelines. A 2025 MDPI Sensors paper~\cite{esp32_mqtt} validated the ESP32 + MQTT architecture for IoT edge-cloud coordination. IoT middleware work by Gazis and Katsiri~\cite{iot_middleware} demonstrated scalable crowd analytics over IoT protocols, while Barthelemy et al.~\cite{barthelemy_2024} validated MQTT-based edge-to-cloud pipelines for real-time violence detection in public transport. An IJRASET 2025 paper~\cite{ijraset_2025} proposed an AI+IoT framework for pre-event capacity prediction with automated alert routing, but relies on cloud ML and lacks real-time flow analysis.

% ===========================================================================
\section{Comparative Analysis with Existing Systems}
% ===========================================================================

Table~\ref{tab:comparison} systematically compares the proposed system against key works across six dimensions: detection methodology, hardware platform, real-time edge capability, stampede risk prediction, automated venue actuation, and deployment venue context.

\begin{table}[h!]
\centering
\caption{Comparison of Proposed System with Existing Literature}
\label{tab:comparison}
\renewcommand{\arraystretch}{1.3}
\small
\begin{tabularx}{\textwidth}{|L{2.8cm}|L{2.5cm}|L{2.0cm}|C{1.2cm}|C{1.4cm}|C{1.4cm}|L{2.2cm}|}
\hline
\rowcolor{tableheadbg}
\color{white}\textbf{Reference} &
\color{white}\textbf{Method} &
\color{white}\textbf{Hardware} &
\color{white}\textbf{Edge RT} &
\color{white}\textbf{Stampede Risk} &
\color{white}\textbf{Actuation} &
\color{white}\textbf{Venue Context} \\
\hline

\rowcolor{lightblue}
\textbf{Proposed System} &
YOLOv8-tiny + Clustering + Flow Vector &
ESP32-CAM + Raspberry Pi (CPU) &
\textbf{Partial} &
\textbf{Proactive} &
\textbf{Yes} &
Stadiums, Stations, Festivals \\
\hline

IEEE 2024~\cite{ieee_yolo_crowd} &
YOLOv3 + Canny / Kalman &
Server GPU &
No &
None &
No &
Generic Surveillance \\
\hline

\rowcolor{altrow}
IEEE 2025~\cite{ieee_thermal} &
Thermal Camera + ML Density &
Dedicated Server &
No &
Reactive &
No &
Generic Indoor \\
\hline

arXiv 2024~\cite{detr_stampede} &
Deformable DETR + Multi-scale Fusion &
GPU Server &
No &
Reactive &
No &
Airports, Stations \\
\hline

\rowcolor{altrow}
Sensors 2024~\cite{iot_middleware} &
IoT WSN + CV on RPi &
Raspberry Pi &
Partial &
None &
No &
Public Spaces \\
\hline

ScienceDirect 2023~\cite{caem_2023} &
YOLOv3 + CAEM Simulation &
Server &
Yes &
Reactive &
Simulation only &
Generic Crowds \\
\hline

\rowcolor{altrow}
Sensors 2024~\cite{barthelemy_2024} &
Action Recognition + MQTT &
Edge Server + AWS &
Yes &
None &
No &
Public Transport \\
\hline

ICE-MoCha~\cite{icemocha} &
RF Signal Analysis + ML &
RF Infrastructure &
Yes &
Partial &
No &
Large Open Events \\
\hline

\end{tabularx}
\end{table}

The comparison clearly establishes that no existing published system simultaneously achieves edge-native inference on affordable commodity IoT hardware, proactive crowd flow-vector-based risk prediction, and automated physical actuation of venue entry/exit signals. The proposed system is the only entry satisfying all three criteria using entirely off-the-shelf, low-cost hardware.

% ===========================================================================
\section{Identified Research Gaps}
% ===========================================================================

Based on the review of over 15 papers from IEEE Xplore, Sensors, Scientific Reports, ScienceDirect, and arXiv (2019--2026), the following three critical research gaps are formally identified.

\subsubsection*{Gap 1: No Integrated Flow-Vector and Density Risk Scoring on Commodity Edge Hardware}
Existing systems either perform server-side density counting or server-side flow analysis in isolation, but none fuse both modalities into a unified proactive risk scoring model deployable on affordable, commodity edge hardware such as the Raspberry Pi. All reviewed fusion approaches require high-performance GPU servers, making them economically unsuitable for wide-scale venue deployment. There is a clear need for a lightweight, integrated pipeline that combines density estimation and flow-vector analysis under real-time constraints on low-cost hardware.

\subsubsection*{Gap 2: Absence of Closed-Loop Automated Venue Actuation}
The reviewed literature uniformly stops at alerting human operators. No published system closes the control loop by autonomously actuating entry gates, exit signals, or crowd access control systems based on computed risk thresholds. Human-in-the-loop delays in crowd emergency scenarios can be catastrophic; an automated, IoT-driven actuation layer is essential yet entirely absent from current research.

% ===========================================================================
\section{Novelty and Positioning of the Proposed System}
% ===========================================================================

Building on the gaps identified in Section~5, the proposed system occupies a clearly distinct position in the research landscape:

\begin{itemize}[leftmargin=1.5em, itemsep=4pt]
  \item \textbf{Commodity hardware edge deployment:} YOLOv8-tiny inference runs directly on a Raspberry Pi 4/5 without any dependency on cloud servers or specialised AI accelerators, making the system practically deployable in real venues at minimal cost.
  \item \textbf{Proactive risk scoring:} The system computes flow vector divergence and density gradient trends to predict crowd crush conditions 30--90 seconds before the critical threshold is breached, unlike reactive detection systems.
  \item \textbf{Closed-loop actuated safety:} The MQTT-to-actuator pipeline enables real-time autonomous adjustment of entry/exit signals via an ESP32 microcontroller — the first such closed-loop crowd safety architecture on commodity IoT hardware in the reviewed literature.
  \item \textbf{Environment-aware spatial risk mapping:} Clustering-based spatial analysis incorporates the physical geometry of the camera's specific field of view, allowing density thresholds to be explicitly tailored to the physical chokepoint or area being monitored.
  \item \textbf{Scalable single-camera architecture:} The current demonstration implements a complete pipeline on a single Raspberry Pi node and camera setup. Because the node operates independently and communicates through a shared MQTT broker, the architecture directly generalises to $n$ cameras deployed across $n$ venue locations with zero changes required to the core software.
  \item \textbf{Cost-effective scalability:} Per-node hardware cost under \$30--\$50 USD (ESP32-CAM + Raspberry Pi) enables multi-node deployment across an entire venue without requiring specialised infrastructure.
\end{itemize}

% ===========================================================================
\section{Methodology}
% ===========================================================================

The proposed system is designed as a fully edge-native, closed-loop crowd safety pipeline. It consists of five tightly integrated stages: (1) image acquisition and hardware configuration, (2) person detection via YOLOv8-tiny, (3) spatial density mapping using clustering, (4) crowd flow vector analysis, and (5) risk scoring with MQTT-based alert and actuation. 

\subsection{Hardware Configuration}

The prototype is implemented as a \textbf{single sensing node setup} consisting of two components: a single ESP32-CAM module for image capture and one Raspberry Pi 4 as the edge compute unit. The ESP32-CAM streams MJPEG video frames over Wi-Fi to the Raspberry Pi, which runs the complete detection, analysis, and risk-scoring pipeline on its CPU. A separate ESP32 microcontroller receives actuation commands via MQTT and drives the physical actuators (servo motor, LEDs, buzzer).

The Raspberry Pi executes the quantised YOLOv8-tiny model using the NCNN or ONNX Runtime inference backend, achieving inference latency in the range of 200--500\,ms per frame at 320$\times$320 resolution — sufficient for the 3--10 FPS operating rate required by the crowd safety pipeline. This frame rate is entirely adequate for detecting crowd density trends and issuing proactive alerts, given the system's 30--90 second predictive alert horizon.

The dashboard runs on a connected laptop, subscribing to the Raspberry Pi's MQTT topics and visualising risk scores, density heatmaps, and alert states in real time.

\textbf{Scalability note:} The current single-camera demonstration successfully showcases the complete pipeline for one venue location. In a real-world deployment, $n$ identical independent nodes — each containing its own Raspberry Pi and camera unit — can be distributed across $n$ venue locations (entry gates, corridors, concourses). Each node simply publishes to its own MQTT topic (e.g., \texttt{venue/gate-1/risk}, \texttt{venue/corridor-2/risk}), and the central dashboard aggregates all feeds simultaneously. No software changes are required to scale from a single camera to $n$ cameras; only additional hardware units need to be deployed.

\subsection{Person Detection with YOLOv8-tiny}

The YOLOv8-tiny model (Ultralytics, 2023) is selected as the detection backbone due to its optimal speed-accuracy trade-off on resource-constrained hardware~\cite{yolov8_edge_2026}. The model is pre-trained on the COCO dataset and fine-tuned on crowd-specific datasets — ShanghaiTech Part B, JHU-Crowd++, and a custom venue footage dataset — to improve recall in high-density, partially occluded scenes.

The input pipeline captures frames at 3--10\,FPS depending on scene complexity. Each frame is pre-processed (resize to 320$\times$320, normalise) and passed through the YOLOv8-tiny inference engine on the Raspberry Pi CPU. Detected person bounding boxes with confidence scores above a threshold (default: 0.45) are extracted, with centroid coordinates computed for each detection.

\subsection{Spatial Density Mapping via Clustering}

Detected person centroids from each frame are fed into a spatial clustering pipeline to generate density estimates for the monitored area:

\begin{enumerate}[leftmargin=1.5em, itemsep=3pt]
  \item \textbf{Field of View Calibration:} The camera frame is calibrated to the physical space it monitors. The entire field of view acts as the primary region of interest and is assigned a configurable maximum safe occupancy threshold based on its specific real-world dimensions (e.g., a specific corridor or gate).
  \item \textbf{Density clustering:} DBSCAN (Density-Based Spatial Clustering of Applications with Noise) is applied to centroid coordinates to identify high-density clusters and isolate noise (lone individuals). Cluster parameters ($\varepsilon$, minPts) are calibrated for the monitored area based on expected crowd geometry.
  \item \textbf{Heatmap generation:} A 2D Gaussian kernel is applied at each detection centroid to produce a continuous density heatmap overlaid on the camera frame. Average density ($\rho$, persons/m$^2$) is computed from the heatmap for the monitored area.
\end{enumerate}

Area density is classified into four risk levels: \textit{Safe} ($\rho < 2.0$\,p/m$^2$), \textit{Caution} ($2.0 \leq \rho < 3.5$\,p/m$^2$), \textit{Warning} ($3.5 \leq \rho < 5.0$\,p/m$^2$), and \textit{Critical} ($\rho \geq 5.0$\,p/m$^2$) — thresholds informed by the Fruin Level-of-Service model used in crowd safety engineering.

\subsection{Optical Flow Vector Analysis}

Crowd flow directionality is computed using the Farneback dense optical flow algorithm applied between consecutive frames~\cite{farneback_2025}. For the monitored area, the following flow metrics are extracted:

\begin{itemize}[leftmargin=1.5em, itemsep=3pt]
  \item \textbf{Mean flow magnitude} ($|\bar{v}|$): Average speed of crowd movement in the area, expressed in pixels/frame and converted to m/s using camera calibration.
  \item \textbf{Flow divergence} ($\nabla \cdot \vec{v}$): Measures whether flow vectors are converging (negative divergence — compression, high crush risk) or dispersing (positive divergence — safe egress).
  \item \textbf{Flow uniformity} ($\sigma_\theta$): Standard deviation of flow vector angles. Low uniformity indicates consistent directional movement; high uniformity indicates turbulent, multi-directional crowd behaviour — a known precursor to panic events.
  \item \textbf{Counter-flow detection:} Situations where opposing flow vectors are detected (bidirectional movement) are flagged with elevated risk, as counter-flow is a primary mechanical cause of crowd crush.
\end{itemize}

\subsection{Risk Score Computation and Predictive Alerting}

A composite risk score $R$ is computed per frame for the monitored area as a weighted combination of density and flow metrics:

\[
R = w_1 \cdot \hat{\rho} + w_2 \cdot |\hat{v}| + w_3 \cdot (1 - \widehat{\nabla \cdot v}) + w_4 \cdot \hat{\sigma}_\theta
\]

where $\hat{(\cdot)}$ denotes min-max normalised values, and weights $w_1 = 0.40$, $w_2 = 0.20$, $w_3 = 0.25$, $w_4 = 0.15$ are empirically calibrated. $R \in [0, 1]$, with values above 0.75 triggering \textit{Critical Alert} and values above 0.55 triggering \textit{Warning Alert}.

A short-term risk trend predictor uses a sliding window of the past 10 risk scores to fit a linear regression and extrapolate 30--90 seconds ahead, enabling proactive alert issuance before the critical threshold is physically reached. The lower frame rate of CPU-based inference is compensated by this longer prediction horizon, maintaining effective system responsiveness.

\subsection{MQTT Communication and Dashboard}

Risk scores, density maps, flow vectors, and alert states are serialised as JSON payloads and published by the Raspberry Pi to an MQTT broker (Mosquitto) running on the connected laptop. The central dashboard — a web-based React application — subscribes to the node's topic and renders:

\begin{itemize}[leftmargin=1.5em, itemsep=3pt]
  \item Live density heatmap overlay on the camera frame
  \item Risk score gauge with historical trend line
  \item Flow vector visualisation
  \item Alert log with timestamp, location label, risk level, and triggered actuation event
\end{itemize}

In a multi-camera deployment context, the dashboard subscribes to all $n$ node topics simultaneously and displays each location's status independently, requiring no architectural changes beyond adding additional MQTT topic subscriptions.

\subsection{Automated Entry/Exit Signal Actuation}

Upon exceeding configurable risk thresholds, the dashboard backend publishes actuation commands via MQTT to an ESP32 microcontroller. To physically demonstrate closed-loop control in our prototype, we utilize standard, cost-effective actuator components to simulate entry and exit barriers:

\begin{itemize}[leftmargin=1.5em, itemsep=3pt]
  \item \textbf{Entry Gate (Open/Close Simulation):} A standard Servo Motor (SG90 or MG995) is used. By attaching a cardboard or acrylic arm, the servo physically rotates to simulate a barrier gate opening and closing. This is controlled directly from the ESP32 via PWM and provides a highly visible, low-cost (approx. ₹100–200) mechanical demonstration of gate actuation.
  \item \textbf{Entry Halt / Warning Indicators:} A simple but universally understood Red/Green LED indicator system is deployed (Red = halt, Green = safe), allowing for immediate comprehension by any examiner.
  \item \textbf{Audible Alerts:} A standard piezoelectric buzzer (approx. ₹30–50) is connected to trigger an immediate, audible alarm when the calculated risk crosses the Critical threshold.
\end{itemize}

Using these physical components, the single-camera prototype demonstrates the following specific actuation behaviours based on risk levels:

\begin{itemize}[leftmargin=1.5em, itemsep=3pt]
  \item \textbf{Entry throttling (Warning level):} The servo motor rotates the barrier arm to a half-open position, simulating a single-person gate mode. An amber/red LED combination activates to signal restricted entry.
  \item \textbf{Entry halt (Critical level):} The servo rotates to the fully closed position, simulating complete gate closure. A red LED activates and the buzzer sounds a continuous alert tone.
  \item \textbf{Safe state:} The servo returns to the fully open position and the green LED activates when the risk score drops below the Warning threshold.
\end{itemize}

These actuator behaviours simulate the physical responses that would be deployed at real venue gates in a full-scale system. The MQTT-based command interface is identical whether driving a prototype servo or a real electromechanical gate controller, ensuring the demonstration is a faithful representation of real-world operation.

The full pipeline — from frame capture on the ESP32-CAM to actuation command received by the ESP32 — is designed to complete within 500\,ms end-to-end under CPU-based inference, well within the safety-relevant response window given the system's 30--90 second predictive alert horizon.

\subsection{Training and Dataset Strategy}

The YOLOv8-tiny model is fine-tuned using a transfer learning approach:

\begin{enumerate}[leftmargin=1.5em, itemsep=3pt]
  \item \textbf{Base pre-training:} COCO person class weights serve as the starting point.
  \item \textbf{Domain fine-tuning:} ShanghaiTech Part B (dense outdoor crowds), JHU-Crowd++ (diverse event venues), and UCF-CC-50 are used for dense-scene adaptation.
  \item \textbf{Venue-specific augmentation:} Overhead perspective transforms, synthetic occlusion masks, and illumination augmentation (day/evening/night) are applied to simulate real venue camera geometry.
  \item \textbf{Prototype validation:} The single-camera prototype is validated by placing groups of people in front of the camera and verifying that density thresholds, risk scores, and actuation responses trigger correctly and consistently.
\end{enumerate}

% ===========================================================================
\section{Conclusion}
% ===========================================================================

This report has presented the complete rationale, literature grounding, and methodology of the Intelligent Crowd Flow \& Density Safety Oracle. The literature survey spanning over 15 publications (2019--2026) identified three critical research gaps: the absence of integrated flow-vector and density risk scoring on commodity edge hardware, the complete lack of closed-loop automated venue actuation, and the non-existence of environment-aware spatial risk mapping informed by venue geometry in practical low-cost deployments. The proposed system addresses all three gaps through a novel edge-native pipeline combining YOLOv8-tiny detection on Raspberry Pi hardware, DBSCAN-based spatial density mapping, Farneback optical flow analysis, and MQTT-driven physical actuation via ESP32 — all using entirely off-the-shelf, affordable components. The current demonstration is implemented as a single-camera, single-node system showcasing the complete pipeline end-to-end; the decentralised, MQTT-based architecture generalises directly to $n$ camera nodes for real-world venue deployment without any software modifications.

% ===========================================================================
\section*{References}
% ===========================================================================
\addcontentsline{toc}{section}{References}

\begin{thebibliography}{99}

\bibitem{ieee_yolo_crowd}
IEEE Xplore, ``Deep Learning-Based Crowd Surveillance and Density Estimation,''
\textit{IEEE Conference Publication}, 2024.
DOI: 10.1109/[conf].10882436

\bibitem{ieee_thermal}
IEEE Xplore, ``Real-Time Crowd Density Estimation and Stampede Risk Assessment System Using Thermal Camera,''
\textit{IEEE Conference Publication}, 2025.
DOI: 10.1109/[conf].10818074

\bibitem{detr_stampede}
X.~Chen et al., ``Stampede Alert Clustering Algorithmic System Based on Tiny-Scale Strengthened DETR,''
Beijing Jiaotong University, arXiv:2404.10359, 2024.

\bibitem{iot_middleware}
A.~Gazis and E.~Katsiri, ``Streamline Intelligent Crowd Monitoring with IoT Cloud Computing Middleware,''
\textit{Sensors}, vol.~24, no.~11, p.~3643, MDPI, 2024.
DOI: 10.3390/s24113643

\bibitem{caem_2023}
``A machine vision-based method for crowd density estimation and evacuation simulation,''
\textit{Safety Science}, ScienceDirect, 2023.
DOI: 10.1016/j.ssci.2023.002278

\bibitem{barthelemy_2024}
J.~Barthelemy et al., ``Safety After Dark: A Privacy Compliant and Real-Time Edge Computing Intelligent Video Analytics for Safer Public Transportation,''
\textit{Sensors}, vol.~24, p.~8102, 2024.
DOI: 10.3390/s24248102

\bibitem{icemocha}
``ICE-MoCha: Intelligent Crowd Engineering using Mobility Characterization and Analytics,''
\textit{PMC / IEEE}, 2019.

\bibitem{esp32_mqtt}
``Design and Implementation of ESP32-Based Edge Computing for Object Detection,''
\textit{MDPI Sensors}, vol.~25, no.~6, p.~1656, 2025.

\bibitem{ieee_xai_crowd}
S.~R.~Alotaibi et al., ``Integrating Explainable Artificial Intelligence with Advanced Deep Learning Model for Crowd Density Estimation in Real-World Surveillance Systems,''
\textit{IEEE Access}, 2025.

\bibitem{yolov8_edge_2026}
``Review of Large YOLOv8 and RT-DETR Energy Efficiency on Edge Devices for Real-Time Detection,''
\textit{Scientific Reports}, 2026.
DOI: 10.1038/s41598-026-46453-6

\bibitem{survey2025}
``A Survey of Deep Learning Methods for Density Estimation and Crowd Counting,''
\textit{Vicinagearth}, Springer, 2025.
DOI: 10.1007/s44336-024-00011-8

\bibitem{mmflux}
A.~Pallaprolu et al., ``mmFlux: Crowd Flow Analytics with Commodity mmWave MIMO Radar,''
\textit{npj Wireless Technology}, 2025.

\bibitem{farneback_2025}
``Density Estimation and Crowd Counting (Event-Driven Farneback Sampling),''
arXiv:2511.09723, 2025.

\bibitem{nature_crowd_cnn}
``Deep Convolutional Neural Network-Based Enhanced Crowd Density Monitoring for Intelligent Urban Planning on Smart Cities,''
\textit{Scientific Reports}, 2025.
DOI: 10.1038/s41598-025-90430-4

\bibitem{langevin}
``Understanding Crowd Flow Movements Using Active-Langevin Model,''
arXiv:2003.05626, 2020.

\bibitem{ijraset_2025}
``AI and IoT-Based Crowd Capacity Prediction and Event Safety Management System Using Historical Event Analysis,''
\textit{IJRASET}, 2025.
DOI: 10.22214/ijraset.2025.76375

\bibitem{yolo}
J.~Redmon et al., ``You Only Look Once: Unified, Real-Time Object Detection,''
\textit{IEEE CVPR}, 2016.

\end{thebibliography}

\end{document}
